% ===========================================================================
\section{Codes, length, and decodability}\label{sec:codes}
\warmth{0}\quad\whisper{A codeword names one symbol; an extended code must preserve a whole sequence.}

\begin{strand}
Coding begins with a map from source symbols to finite strings. Short strings
save space, but distinct symbols and distinct source sequences must remain
distinguishable. This is why expected length and decodability have to be studied
together.
\end{strand}

\block{From a source symbol to a codeword}

\begin{definition}[source code]\label{def:source-code}
Let $X$ be a random variable with finite alphabet $\X$. A $D$-ary source code is
a map
\[
  C:\X\longrightarrow\D^*,
\]
where $\D^*$ is the set of \answerin[3.9cm]{finite-length strings} over a
$D$-ary alphabet $\D$. The string $C(x)$ is the \keyword{codeword} of $x$, and
$|C(x)|$ is its length.
\studentrecall{The blank is ``finite-length strings.''}
\end{definition}

For instance, if $\X=\{\text{red},\text{blue}\}$ and $\D=\{0,1\}$, one may set
$C(\text{red})=00$ and $C(\text{blue})=11$.

\begin{definition}[expected length]\label{def:expected-length}
If $X$ has probability mass function $p$, the expected length of $C$ is
\[
  L(C)=\E[|C(X)|]
      =\sum_{x\in\X}p(x)|C(x)|.
\]
\end{definition}

\begin{example}[an entropy-matched binary code]\label{ex:dyadic-code}
Let $\X=\{1,2,3,4\}$, with probabilities and codewords
\[
\begin{array}{c|cccc}
x&1&2&3&4\\ \hline
p(x)&\tfrac12&\tfrac14&\tfrac18&\tfrac18\\
C(x)&0&10&110&111
\end{array}
\]
The codeword lengths are $1,2,3,3$, so
\[
  L(C)=\tfrac12\cdot1+\tfrac14\cdot2
       +\tfrac18\cdot3+\tfrac18\cdot3
       =\answerin[1.4cm]{\tfrac74}\ \text{bits}.
\]
Since $|C(x)|=-\log_2 p(x)$ for every symbol here,
\[
  \Ent(X)=\sum_xp(x)\log_2\frac1{p(x)}
         =\answerin[1.4cm]{\tfrac74}\ \text{bits}=L(C).
\]
\studentrecall{Both blanks are $\tfrac74$ bits.}
\end{example}

\begin{yourturn}
Why does the equality $L(C)=\Ent(X)$ require no averaging trick in this example?
\studentrecall{Here $|C(x)|=-\log_2p(x)$ for every $x$, so equality holds term by term.}
\ifshowanswers\else\workspace[2]\fi
\answerprose{The equality holds term by term: each length satisfies
$|C(x)|=-\log_2p(x)$. Multiplying by $p(x)$ and summing therefore gives
$L(C)=\Ent(X)$.}
\end{yourturn}

\block{Three levels of decodability}

\begin{definition}[nonsingular code]\label{def:nonsingular}
A code $C$ is \keyword{nonsingular} if different source symbols have different
codewords:
\[
  x\ne x'\quad\Longrightarrow\quad C(x)\ne C(x').
\]
Equivalently, $C$ is \answerin[2.2cm]{injective}.
\studentrecall{The missing word is ``injective.''}
\end{definition}

Nonsingularity only protects one symbol at a time. To encode a sequence, extend
$C$ by concatenation:
\[
  C^*:\X^*\longrightarrow\D^*,\qquad
  C^*(x_1x_2\cdots x_n)=C(x_1)C(x_2)\cdots C(x_n).
\]

\begin{definition}[uniquely decodable code]\label{def:ud}
A code is \keyword{uniquely decodable} if its extension $C^*$ is nonsingular.
Thus every encoded string has at most one source sequence that produced it.
\end{definition}

\begin{definition}[prefix code]\label{def:prefix}
A code is a \keyword{prefix code}, or instantaneous code, if no codeword is a
prefix of any other codeword.
\end{definition}

Every prefix code is uniquely decodable, and every uniquely decodable code is
nonsingular:
\[
  \{\text{prefix codes}\}
  \subset
  \{\text{uniquely decodable codes}\}
  \subset
  \{\text{nonsingular codes}\}.
\]
The converses need not hold.

\begin{example}[uniquely decodable but not prefix]\label{ex:ud-not-prefix}
For $\X=\{A,B,C,D\}$, consider
\[
  C(A)=0,\qquad C(B)=01,\qquad C(C)=011,\qquad C(D)=111.
\]
This is not a prefix code because $0$ is a prefix of both $01$ and $011$.
Nevertheless, it is uniquely decodable. Removing an initial $0$ from a longer
codeword leaves only $1$ or $11$; comparing these suffixes with $111$ swaps
$1\leftrightarrow11$ and never produces the empty suffix. Hence no two source
sequences have the same concatenation.
\end{example}

\begin{yourturn}
The code $\{0,01,10\}$ is nonsingular but not uniquely decodable. Exhibit two
different source sequences with the same encoded string.
\studentrecall{$AC$ and $BA$ both encode to $010$.}
\ifshowanswers\else\workspace[1]\fi
\answerprose{$AC$ and $BA$ both encode as $010$.}
\end{yourturn}

\vspace{1.2cm}
\recall{Why does nonsingularity not guarantee decodability after concatenation?}
